\section{Preliminaries}
\label{sec:background}

\subsection{Gradient flow and the scalar auxiliary variable}

Let \(F:\R^d\to\R\) be differentiable and bounded from below. Gradient descent
is the forward Euler discretization of the Euclidean gradient flow
\(\dot\theta=-\nabla F(\theta)\), along which
\(\tfrac{\mathrm d}{\mathrm dt}F=-\norm{\nabla F}^{2}\leq0\). Choose \(C\)
with \(F(\theta)+C>0\) on the states of interest and set
\(Q(\theta)=\sqrt{F(\theta)+C}\) and \(g(\theta)=\nabla F(\theta)\). In the SAV
reformulation for gradient flows \cite{shen2018sav,shen2019savreview}, an
auxiliary scalar \(q\) tracks \(Q(\theta(t))\) and the modified energy
\(q^2\) satisfies
\(\tfrac{\mathrm d}{\mathrm dt}q^2=-\norm{\dot\theta}^{2}\leq0\).
Convergence and error analysis for the resulting schemes
\cite{shen2018convergence} and a related energy-stable variant
\cite{huang2020newsav} are available.

For the discrete formulation, let \(g^n=g(\theta^n)\) and
\(Q^n=Q(\theta^n)\), initialize \(q^0=Q(\theta^0)\), and let \(\eta>0\) be
the step size. The first-order SAV update is
\begin{subequations}
\label{eq:discrete-sav-coupled}
\begin{align}
  \frac{\theta^{n+1}-\theta^n}{\eta}
  &=-\frac{q^{n+1}}{Q^n}g^n,
  \label{eq:discrete-sav-theta}\\
  q^{n+1}-q^n
  &=\frac{(g^n)^{\mathsf T}(\theta^{n+1}-\theta^n)}{2Q^n}.
  \label{eq:discrete-sav-q}
\end{align}
\end{subequations}
Writing \(\Delta\theta^{n+1}=\theta^{n+1}-\theta^n\) and eliminating
\(q^{n+1}\) gives
\begin{equation}
  \left(\eta^{-1}I_d+B_1^n\right)\Delta\theta^{n+1}
  =-\frac{q^n}{Q^n}g^n,
  \qquad
  B_1^n=
  \frac{g^n(g^n)^{\mathsf T}}{2(Q^n)^2}.
  \label{eq:sav-eliminated}
\end{equation}
Here \(I_d\) is the \(d\times d\) identity matrix, and identities of other
sizes carry the corresponding subscript. The correction \(B_1^n\) has rank
one, and \(q^n/Q^n=1\) when the stored scalar matches \(Q(\theta^n)\).

\subsection{PB--SAV for composite objectives}

Consider an additive objective
\begin{equation}
  F(\theta)=\sum_{i=1}^{m}E_i(\theta),
  \qquad
  E_i(\theta)+C_i>0,
  \qquad
  C=\sum_{i=1}^{m}C_i.
  \label{eq:component-decomposition}
\end{equation}
Multiple-SAV methods assign a separate auxiliary variable to each energy
component \cite{cheng2018msav}. PB--SAV instead retains one auxiliary variable
for the total shifted objective and builds the correction from the individual
components \cite{zhang2026pbsav}.

At \(\theta^n\), define
\begin{equation}
  Q_i^n=\sqrt{E_i(\theta^n)+C_i},
  \qquad
  g_i^n=\nabla E_i(\theta^n),
  \qquad
  g^n=\sum_{i=1}^{m}g_i^n,
  \label{eq:component-notation-preview}
\end{equation}
so that \((Q^n)^2=\sum_{i=1}^{m}(Q_i^n)^2\). The component matrix is
\begin{equation}
  V^n=
  \begin{bmatrix}
  \dfrac{g_1^n}{\sqrt{2}Q_1^n}
  &\cdots&
  \dfrac{g_m^n}{\sqrt{2}Q_m^n}
  \end{bmatrix},
  \qquad
  B_m^n=V^n(V^n)^{\mathsf T}.
  \label{eq:pbsav-preview}
\end{equation}

\subsection{Component curvature}
\label{sec:recalled-geometry}

The PB--SAV pullback is a Gauss--Newton matrix for the transformed residuals
\cite{zhang2026pbsav}
\begin{equation}
  \psi_i(\theta)=\sqrt{2(E_i(\theta)+C_i)},
\end{equation}
with \(\psi=(\psi_1,\ldots,\psi_m)^{\mathsf T}\) and
\(F+C=\tfrac12\norm{\psi}^2\). If \(J_\psi(\theta)\) denotes its Jacobian,
differentiation gives
\begin{equation}
  B_m^n=J_\psi(\theta^n)^{\mathsf T}J_\psi(\theta^n).
  \label{eq:background-pullback-gn}
\end{equation}
For scalar least-squares components \(E_i=\tfrac12 r_i^2\) with \(C_i>0\),
this formula becomes
\begin{equation}
  B_m
  =\sum_{i=1}^m
  \frac{r_i^2}{r_i^2+2C_i}\,
  \nabla r_i\nabla r_i^{\mathsf T}.
  \label{eq:residual-pullback-curvature}
\end{equation}
Thus the component pullback approximates the Gauss--Newton matrix
\(G=\sum_i\nabla r_i\nabla r_i^{\mathsf T}\) when the shifts are small
relative to the squared residuals. For twice continuously differentiable
residuals,
\begin{equation}
  \nabla^2F=G+\sum_i r_i\nabla^2r_i,
  \label{eq:least-squares-hessian}
\end{equation}
so the residual-weighted second derivatives determine the difference between
Gauss--Newton and Hessian curvature \cite{bjorck1996least}.

The relation to the aggregate correction follows from a weighted variance
identity \cite{zhang2026pbsav}. For any \(z\in\R^d\), define
\begin{equation}
  y_i^n=\frac{z^{\mathsf T}g_i^n}{(Q_i^n)^2},
  \qquad
  \bar y^n=\frac{\sum_i(Q_i^n)^2y_i^n}{\sum_i(Q_i^n)^2}.
\end{equation}
Then
\begin{equation}
  z^{\mathsf T}(B_m^n-B_1^n)z
  =\frac12\sum_i(Q_i^n)^2(y_i^n-\bar y^n)^2\geq0.
  \label{eq:background-variance-gap}
\end{equation}
Thus \(B_m^n-B_1^n\succeq0\).
